% ============================================================
%  IACV 2025/26 Project — Synthetic Pose Estimation Pipeline
%  Mathematical Foundations of the Full Pipeline (steps 1-7)
%  Ready to compile on Overleaf (pdfLaTeX).
% ============================================================
\documentclass[11pt]{article}

\usepackage[utf8]{inputenc}
\usepackage[T1]{fontenc}
\usepackage{amsmath, amssymb, amsthm, mathtools}
\usepackage{bm}
\usepackage[margin=2.4cm]{geometry}
\usepackage{xcolor}
\usepackage{hyperref}
\usepackage{enumitem}

\hypersetup{colorlinks=true, linkcolor=blue!60!black, urlcolor=blue!60!black}

% --- Notation shortcuts ---
\newcommand{\R}{\mathbb{R}}
\newcommand{\SO}[1]{\mathrm{SO}(#1)}
\newcommand{\so}[1]{\mathfrak{so}(#1)}
\newcommand{\SE}[1]{\mathrm{SE}(#1)}
\newcommand{\norm}[1]{\left\lVert #1 \right\rVert}
\newcommand{\skewmat}[1]{\left[ #1 \right]_{\times}}
\newcommand{\Rgt}{\mathbf{R}_{\text{gt}}}
\newcommand{\tgt}{\mathbf{t}_{\text{gt}}}
\newcommand{\wgt}{\boldsymbol{\omega}_{\text{gt}}}
\newcommand{\Pcam}{\mathbf{P}_{\text{cam}}}
\newcommand{\Pmark}{\mathbf{P}_{\text{mk}}}
\newcommand{\KK}{\mathbf{K}}
\newcommand{\Hom}{\mathbf{H}}

\title{Synthetic Pose Estimation Pipeline \\[0.3em]
       \large Mathematical Foundations --- Steps 1 to 7}
\author{IACV 2025/26 Project}
\date{\today}

\begin{document}
\maketitle
\tableofcontents
\vspace{1em}

% ============================================================
\section{Overview and Notation}
% ============================================================

We synthesize the image a calibrated pinhole camera would observe of a
known planar square marker placed at a known pose, then re-estimate the
pose from the synthesized image points using our own implementation of
\verb|SOLVEPNP_IPPE_SQUARE|. The full pipeline:

\begin{enumerate}[label=\textbf{Step \arabic*.}, leftmargin=*]
  \item Generate a 3D plane (the marker plane).
  \item Define a square on that plane of side $L$.
  \item Define a camera with known intrinsics $\KK$.
  \item Choose a known marker pose $(\Rgt, \tgt)$ relative to the camera.
  \item Project the four 3D corners into the image.
  \item Re-estimate the pose with our IPPE\_SQUARE (and compare with OpenCV).
  \item Compare the estimated pose with the ground truth.
\end{enumerate}

Steps 1--5 are pure forward simulation (no estimation, no noise). Step 6 is
the actual pose-estimation algorithm. Step 7 quantifies how well the
estimator recovered the chosen ground truth. The two experiments asked for
in the brief (two-pose ambiguity sweep; Jacobian conditioning study) are
\emph{after} this pipeline and are not covered in this document.

\subsection*{Notation}
\begin{itemize}[leftmargin=*]
  \item Vectors in bold lowercase: $\mathbf{t}, \boldsymbol{\omega}$. Matrices in bold uppercase: $\mathbf{R}, \KK$.
  \item Points carry a subscript naming their frame:
        $\Pmark$ is in the marker frame, $\Pcam$ in the camera frame.
  \item Pixel coordinates: $\mathbf{p} = (u, v)^\top$.
  \item Coordinate frames are right-handed: in the camera frame, $X$ points right,
        $Y$ points down, $Z$ points into the scene.
\end{itemize}

% ============================================================
\section{Step 1 --- The 3D Marker Plane}
% ============================================================

The marker is a flat object, so the four corners are coplanar. We attach a
coordinate frame to the marker (the \emph{marker frame}) with the plane
chosen as $Z = 0$:
\begin{equation}
  \text{marker plane} \;=\; \{\, (X, Y, 0) \,:\, X, Y \in \R \,\}.
\end{equation}
The origin is placed at the geometric center of the future square. The
$X$ axis points right, $Y$ points up (in the plane), and $Z$ is the outward
normal to the marker surface.

\paragraph{Why $Z = 0$ matters.}
A general projection of a 3D point requires the full $3\times 3$ rotation
matrix. If the third coordinate is identically zero, the third column of
$\mathbf{R}$ multiplies $0$ and drops out:
\begin{equation}
  \mathbf{R} \begin{pmatrix} X \\ Y \\ 0 \end{pmatrix}
  \;=\; X \mathbf{r}_1 + Y \mathbf{r}_2,
\end{equation}
where $\mathbf{r}_1, \mathbf{r}_2$ are the first two columns of $\mathbf{R}$. As we
will see in Step~6, this collapse is the structural reason a single
$3\times 3$ matrix --- a \emph{homography} --- captures all the geometry.

% ============================================================
\section{Step 2 --- The Square on the Plane}
% ============================================================

On the plane $Z = 0$ we place a square of side $L$, centered at the marker
origin and axis-aligned to the marker frame. The four corner coordinates,
in the order used by OpenCV's \verb|SOLVEPNP_IPPE_SQUARE|, are
\begin{equation}
  \Pmark^{(0)} = \begin{pmatrix} -L/2 \\ +L/2 \\ 0 \end{pmatrix}, \;
  \Pmark^{(1)} = \begin{pmatrix} +L/2 \\ +L/2 \\ 0 \end{pmatrix}, \;
  \Pmark^{(2)} = \begin{pmatrix} +L/2 \\ -L/2 \\ 0 \end{pmatrix}, \;
  \Pmark^{(3)} = \begin{pmatrix} -L/2 \\ -L/2 \\ 0 \end{pmatrix}.
\end{equation}
The parenthesized superscript is the corner index; the subscript ``mk''
indicates the marker frame.

\paragraph{Corner ordering matters.}
Pose estimators consume \emph{pairs} of corresponding 3D and 2D corners.
The order of the 3D list must match the order of the 2D observations,
otherwise the correspondences are scrambled and the algorithm finds the
wrong pose (or fails). The order above (counter-clockwise starting at
top-left, viewed from $+Z$) is the convention shared by OpenCV's ArUco
detector and \verb|SOLVEPNP_IPPE_SQUARE|.

% ============================================================
\section{Step 3 --- The Pinhole Camera}
% ============================================================

\subsection{Geometric model}

A pinhole camera maps points in the camera frame to points on a 2D image.
The camera origin is at the world origin; the image plane sits at $Z = f$,
where $f$ is the focal length. By similar triangles, a 3D point
$\Pcam = (X, Y, Z)^\top$ projects to
\begin{equation}
  x = f\cdot\frac{X}{Z}, \qquad y = f\cdot\frac{Y}{Z}.
  \label{eq:similar-triangles}
\end{equation}
The division by $Z$ is the \emph{perspective divide}: it is the only
non-linearity in the model and is responsible for foreshortening (farther
objects look smaller).

\subsection{The intrinsic matrix \texorpdfstring{$\KK$}{K}}

Real images are measured in pixels. The focal length is encoded as
$f_x, f_y$ (in pixels, possibly different along the two axes if the sensor
elements are not square), and the optical axis hits the image at the
\emph{principal point} $(c_x, c_y)$. The pixel coordinates are
\begin{equation}
  u = f_x \cdot \frac{X}{Z} + c_x, \qquad
  v = f_y \cdot \frac{Y}{Z} + c_y.
  \label{eq:pixel-projection}
\end{equation}
Collected into a matrix:
\begin{equation}
  \boxed{\quad
    \KK \;=\; \begin{pmatrix} f_x & 0 & c_x \\ 0 & f_y & c_y \\ 0 & 0 & 1 \end{pmatrix}
  \quad}
\end{equation}
so that the projection has the compact homogeneous form
\begin{equation}
  \tilde{\mathbf{p}} \;=\; \KK\,\Pcam,
  \qquad
  \mathbf{p} \;=\; \begin{pmatrix} \tilde p_1/\tilde p_3 \\ \tilde p_2/\tilde p_3 \end{pmatrix}.
\end{equation}
First line: a linear matrix product. Second line: the perspective divide.

\subsection{Normalized image coordinates}

Apply $\KK^{-1}$ to a pixel $\mathbf{p}$:
\begin{equation}
  \begin{pmatrix} x_n \\ y_n \\ 1 \end{pmatrix}
  \;=\; \KK^{-1}
  \begin{pmatrix} u \\ v \\ 1 \end{pmatrix}
  \;=\;
  \begin{pmatrix} (u - c_x)/f_x \\ (v - c_y)/f_y \\ 1 \end{pmatrix}.
\end{equation}
The pair $(x_n, y_n)$ are the \emph{normalized image coordinates}: the
location where the ray through pixel $\mathbf{p}$ crosses the plane
$Z = 1$ in camera space. They are dimensionless (no pixels, no focal
length) and encode pure ray directions. Pose-estimation algorithms (Step~6)
work in this representation so that the camera-specific quantities $\KK$
do not contaminate the geometric reasoning.

% ============================================================
\section{Step 4 --- Ground-Truth Pose}
% ============================================================

A point in the marker frame is brought into the camera frame by a
rigid-body transform:
\begin{equation}
  \boxed{\quad \Pcam \;=\; \mathbf{R}\,\Pmark + \mathbf{t} \quad}
  \label{eq:rigid-transform}
\end{equation}
with $\mathbf{R} \in \SO{3}$ and $\mathbf{t} \in \R^3$. The pair
$(\mathbf{R}, \mathbf{t})$ is an element of the Special Euclidean group
$\SE{3}$, and is what we mean by \emph{pose}.

\subsection{The rotation group \texorpdfstring{$\SO{3}$}{SO(3)}}

A $3\times 3$ matrix $\mathbf{R}$ is a rotation if and only if
\begin{equation}
  \mathbf{R}^\top \mathbf{R} \;=\; \mathbf{I},
  \qquad
  \det(\mathbf{R}) \;=\; +1.
  \label{eq:so3-constraints}
\end{equation}
The first condition (orthonormal columns) preserves lengths and angles; the
second excludes reflections, which would flip a right-handed frame into a
left-handed one.

\paragraph{Degrees of freedom.}
A $3\times 3$ matrix has $9$ entries. The orthogonality condition
$\mathbf{R}^\top \mathbf{R} = \mathbf{I}$ imposes $6$ independent constraints. That
leaves \textbf{3 DOF} --- the minimum number of parameters needed to
describe a rotation in 3D.

\subsection{Representation 1: Euler angles}

Decompose any rotation as three elementary rotations around the coordinate
axes. The \emph{intrinsic xyz} convention used in our code is
\begin{equation}
  \mathbf{R}(\alpha,\beta,\gamma) \;=\; \mathbf{R}_x(\alpha)\,\mathbf{R}_y(\beta)\,\mathbf{R}_z(\gamma),
\end{equation}
where each $\mathbf{R}_\bullet(\cdot)$ is a $3\times 3$ planar rotation, e.g.
\begin{equation}
  \mathbf{R}_x(\alpha) =
  \begin{pmatrix} 1 & 0 & 0 \\ 0 & \cos\alpha & -\sin\alpha \\ 0 & \sin\alpha & \cos\alpha \end{pmatrix}.
\end{equation}
\textbf{Pros:} human-intuitive --- $(15^\circ, 10^\circ, 5^\circ)$ is easy to picture.\\
\textbf{Cons:} suffer from \emph{gimbal lock} (a configuration where two
axes align and one DOF is lost) and from non-unique decompositions. We use
Euler angles only to \emph{specify} the ground-truth pose conveniently; we
never use them for optimization.

\subsection{Representation 2: Axis--angle (rotation vector)}

Euler's rotation theorem states that every $\mathbf{R} \in \SO{3}$ is a
rotation by some angle $\theta$ around some unit axis $\hat{\mathbf{n}}$. Combine
these into a single 3-vector
\begin{equation}
  \boxed{\quad
    \boldsymbol{\omega} \;=\; \theta\,\hat{\mathbf{n}} \;\in\; \R^3
  \quad}
  \label{eq:axis-angle}
\end{equation}
called the \emph{rotation vector} or \emph{axis--angle vector}. Its
direction is the rotation axis, its magnitude is the rotation angle:
\begin{equation}
  \theta = \norm{\boldsymbol{\omega}}, \qquad
  \hat{\mathbf{n}} = \boldsymbol{\omega}/\norm{\boldsymbol{\omega}}.
\end{equation}
This minimal 3-parameter form is the \textbf{Lie algebra} $\so{3}$ of the
rotation group $\SO{3}$.

\subsection{The Rodrigues formula (exponential map)}

The map from a rotation vector $\boldsymbol{\omega}$ to a rotation matrix
$\mathbf{R}$ is given by the \emph{Rodrigues formula}:
\begin{equation}
  \boxed{\quad
    \mathbf{R}(\boldsymbol{\omega}) \;=\; \mathbf{I}
      + \sin\theta\,\skewmat{\hat{\mathbf{n}}}
      + (1 - \cos\theta)\,\skewmat{\hat{\mathbf{n}}}^2
  \quad}
  \label{eq:rodrigues}
\end{equation}
where $\theta = \norm{\boldsymbol{\omega}}$, $\hat{\mathbf{n}} = \boldsymbol{\omega}/\theta$, and
$\skewmat{\cdot}$ is the \emph{skew-symmetric matrix} built from a 3-vector:
\begin{equation}
  \skewmat{\mathbf{v}} \;=\;
  \begin{pmatrix}
    0 & -v_3 & v_2 \\
    v_3 & 0 & -v_1 \\
    -v_2 & v_1 & 0
  \end{pmatrix},
  \qquad
  \skewmat{\mathbf{v}}\,\mathbf{u} \;=\; \mathbf{v} \times \mathbf{u}.
\end{equation}

\paragraph{Intuition.} Read \eqref{eq:rodrigues} as ``start from the
identity (no rotation), add a first-order term proportional to
$\sin\theta$, and a second-order correction''. It is the closed form of the
matrix exponential
\[
  \mathbf{R}(\boldsymbol{\omega}) \;=\; \exp\!\big(\skewmat{\boldsymbol{\omega}}\big)
  \;=\; \sum_{k=0}^{\infty} \frac{\skewmat{\boldsymbol{\omega}}^k}{k!},
\]
which collapses to a finite sum because $\skewmat{\hat{\mathbf{n}}}^3 = -\skewmat{\hat{\mathbf{n}}}$.

\paragraph{Inverse map.} Given $\mathbf{R}$, recover the rotation vector via
\begin{equation}
  \theta = \arccos\!\left(\frac{\mathrm{trace}(\mathbf{R}) - 1}{2}\right),
  \quad
  \skewmat{\hat{\mathbf{n}}} = \frac{\mathbf{R} - \mathbf{R}^\top}{2\sin\theta},
  \quad
  \boldsymbol{\omega} = \theta\,\hat{\mathbf{n}}.
\end{equation}
The \texttt{scipy.spatial.transform.Rotation} class exposes this as
\verb|.as_rotvec()|.

\subsection{Our concrete ground-truth pose}

In the code we choose
\begin{align}
  \text{Euler (xyz, deg)} & : \quad (\alpha, \beta, \gamma) = (15,\,10,\,5) \\
  \tgt & : \quad (0.05,\, 0.02,\, 0.50)^\top \;\text{m}
\end{align}
which yields
\[
  \Rgt \approx
  \begin{pmatrix}
    +0.9811 & -0.0394 & +0.1897 \\
    +0.0858 & +0.9662 & -0.2432 \\
    -0.1736 & +0.2549 & +0.9513
  \end{pmatrix},
  \qquad
  \wgt \approx
  \begin{pmatrix} 0.2534 \\ 0.1848 \\ 0.0637 \end{pmatrix} \text{rad}.
\]
Note $\norm{\wgt} \approx 18.33^\circ$, the total rotation angle. It is
\emph{not} $15^\circ + 10^\circ + 5^\circ$ because Euler rotations do not
compose by simple addition.

% ============================================================
\section{Step 5 --- Forward Projection}
% ============================================================

Steps 1--4 fix all the inputs. Step 5 chains them into the observed image
points. For each marker corner $\Pmark^{(i)}$:
\begin{equation}
  \boxed{\quad
    \mathbf{p}^{(i)}
    \;=\;
    \pi\!\Big(\, \KK\,(\,\Rgt \Pmark^{(i)} + \tgt\,) \,\Big)
  \quad}
  \label{eq:full-projection}
\end{equation}
where the projection $\pi:\R^3 \to \R^2$ is the perspective divide
\begin{equation}
  \pi\!\begin{pmatrix} \tilde x \\ \tilde y \\ \tilde z \end{pmatrix}
  \;=\;
  \begin{pmatrix} \tilde x / \tilde z \\ \tilde y / \tilde z \end{pmatrix}.
\end{equation}

\paragraph{In the code, Step 5 is split in two sub-steps:}
\begin{enumerate}[label=(\alph*)]
  \item \textbf{Rigid transform.} $\Pcam^{(i)} \;=\; \Rgt \Pmark^{(i)} + \tgt$.
        Brings each marker corner from the marker frame into the camera frame.
  \item \textbf{Pinhole projection.} $\mathbf{p}^{(i)} \;=\; \pi(\KK\,\Pcam^{(i)})$.
        Maps the camera-frame point to pixel coordinates via
        \eqref{eq:pixel-projection}.
\end{enumerate}

\paragraph{Sanity checks at this point:}
\begin{itemize}[leftmargin=*]
  \item All $\Pcam^{(i)}$ have $Z > 0$ (in front of the camera --- \emph{cheirality}).
  \item All $\mathbf{p}^{(i)}$ fall inside $[0, W) \times [0, H)$ (visible on the sensor).
  \item The four $\mathbf{p}^{(i)}$ form a \emph{convex quadrilateral} that looks
        like a perspective-warped square.
\end{itemize}

% ============================================================
\section{Step 6 --- Pose Estimation with IPPE\_SQUARE}
% ============================================================

We now invert the projection: given $\KK$, the four marker corners
$\Pmark^{(i)}$, and the four image points $\mathbf{p}^{(i)}$, recover
$(\mathbf{R}, \mathbf{t})$. We follow the structure of OpenCV's
\verb|SOLVEPNP_IPPE_SQUARE| (Collins \& Bartoli, IJCV 2014).

\subsection{What IPPE\_SQUARE is}

\textbf{IPPE} = \emph{Infinitesimal Plane-Based Pose Estimation}.
For \emph{planar} markers, the projection from the marker plane to the
image plane is a \textbf{homography} $\Hom$ (a $3\times 3$ matrix). IPPE
exploits planarity to obtain a closed-form pose --- no iterative
optimization. The \verb|_SQUARE| specialization exploits the additional
fact that the marker is a \emph{centered, axis-aligned square}, which
simplifies the homography step further.

\textbf{Two solutions.} The planar pose problem has an inherent
two-fold ambiguity: there exist (in general) two distinct poses
$(\mathbf{R}_1, \mathbf{t}_1)$ and $(\mathbf{R}_2, \mathbf{t}_2)$ that project the same 3D
square to image points indistinguishable up to noise. IPPE returns both,
ranked by reprojection error.

\subsection{The pipeline at a glance}

\begin{center}
\begin{tabular}{|c|l|}
\hline
\textbf{Block} & \textbf{Operation} \\
\hline
A & Normalize image points: $\mathbf{q}^{(i)} = \KK^{-1}\,\mathbf{p}^{(i)}$ \\
B & Estimate homography $\Hom$ from $\{\Pmark^{(i)}_{XY}\} \to \{\mathbf{q}^{(i)}\}$ (DLT) \\
C & Decompose $\Hom$ into two candidate poses $(\mathbf{R}, \mathbf{t})$ \\
D & Project each $\mathbf{R}$ onto $\SO{3}$ (Procrustes via SVD) \\
E & Rank the candidates by reprojection error \\
\hline
\end{tabular}
\end{center}

\subsection{Block A --- Normalize image points}

Apply $\KK^{-1}$ to each observed pixel:
\begin{equation}
  \mathbf{q}^{(i)} \;=\; \KK^{-1}\,\mathbf{p}^{(i)} \;\in\; \R^2.
\end{equation}
This removes $\KK$ from the rest of the pipeline. The homography we
compute next will map the marker plane directly to the $Z = 1$ normalized
plane, without any factor of $\KK$ in it.

\subsection{Block B --- Homography estimation by DLT}

For a planar point $\Pmark = (X, Y, 0)^\top$, the projection collapses to
\begin{equation}
  \tilde{\mathbf{q}} \;=\; [\mathbf{r}_1 \,|\, \mathbf{r}_2 \,|\, \mathbf{t}]
  \begin{pmatrix} X \\ Y \\ 1 \end{pmatrix},
\end{equation}
so the $3\times 3$ matrix $\Hom \;=\; [\mathbf{r}_1\,|\,\mathbf{r}_2\,|\,\mathbf{t}]$
(up to a global scale) is the \emph{homography} from the marker plane to
the normalized image plane. With four correspondences
$(X_i, Y_i) \leftrightarrow (x_n^{(i)}, y_n^{(i)})$ we estimate $\Hom$
linearly via the \textbf{Direct Linear Transform (DLT)}.

\paragraph{From homogeneous equality to linear equations.}
Write the action of $\Hom$ in coordinates:
\begin{equation}
  \begin{pmatrix} \tilde u \\ \tilde v \\ \tilde w \end{pmatrix}
  \;=\; \Hom \begin{pmatrix} X \\ Y \\ 1 \end{pmatrix},
\end{equation}
the observed normalized point is $(x_n, y_n) = (\tilde u/\tilde w,\,\tilde v/\tilde w)$.
Cross-multiplying:
\begin{align}
  x_n\,\tilde w - \tilde u &= 0, \\
  y_n\,\tilde w - \tilde v &= 0.
\end{align}
Substituting $\tilde u, \tilde v, \tilde w$ in terms of the entries of
$\Hom = \{h_{ij}\}$ yields two linear equations \emph{per point} in the
9 unknowns:
\begin{equation}
  \underbrace{
  \begin{pmatrix}
    -X & -Y & -1 &  0 &  0 & 0 & x_n X & x_n Y & x_n \\
     0 &  0 &  0 & -X & -Y & -1 & y_n X & y_n Y & y_n
  \end{pmatrix}}_{2\times 9 \text{ block of } \mathbf{A}_i}
  \;
  \underbrace{
  \begin{pmatrix} h_{11} \\ h_{12} \\ \vdots \\ h_{33} \end{pmatrix}}_{\mathbf{h} \;\in\; \R^9}
  \;=\; \mathbf{0}.
\end{equation}

\paragraph{Stacking four points.}
With $N = 4$ correspondences, the system is $\mathbf{A}\,\mathbf{h} = \mathbf{0}$
with $\mathbf{A} \in \R^{8 \times 9}$. The trivial solution $\mathbf{h} = \mathbf{0}$ is useless;
the meaningful one is found by SVD:
\begin{equation}
  \mathbf{A} \;=\; \mathbf{U}\,\boldsymbol{\Sigma}\,\mathbf{V}^\top,
  \qquad
  \mathbf{h} \;=\; \text{last column of } \mathbf{V} \;=\; \text{last row of } \mathbf{V}^\top.
\end{equation}
This is the right singular vector corresponding to the smallest singular
value, the choice that minimizes $\norm{\mathbf{A}\mathbf{h}}^2$ subject to
$\norm{\mathbf{h}} = 1$ (Eckart--Young theorem). Reshape $\mathbf{h}$ into a
$3 \times 3$ matrix to obtain $\Hom$.

\subsection{Block C --- Decomposing \texorpdfstring{$\Hom$}{H} into pose}

Because we normalized the image points in Block~A, the homography we
computed satisfies
\begin{equation}
  \Hom \;=\; [\mathbf{r}_1 \,|\, \mathbf{r}_2 \,|\, \mathbf{t}] \qquad \text{(up to a global scale } s\text{)}.
\end{equation}
A rotation matrix has unit-norm columns, so the scale is fixed by
$\norm{\mathbf{r}_1} = 1$:
\begin{equation}
  s \;=\; \norm{\Hom_{:,1}}, \qquad
  \mathbf{r}_1 \;=\; \Hom_{:,1}/s, \quad
  \mathbf{r}_2 \;=\; \Hom_{:,2}/s, \quad
  \mathbf{t} \;=\; \Hom_{:,3}/s.
\end{equation}
The third rotation column is recovered via the cross product
\begin{equation}
  \mathbf{r}_3 \;=\; \mathbf{r}_1 \times \mathbf{r}_2,
\end{equation}
which automatically yields a right-handed orthogonal axis.

\paragraph{The sign ambiguity (two solutions).}
A homography is defined up to a non-zero scalar: $\Hom$ and $-\Hom$
represent the same map between projective planes. Therefore both
$+s$ and $-s$ produce valid decompositions, giving two candidate poses
$(\mathbf{R}^+, \mathbf{t}^+)$ and $(\mathbf{R}^-, \mathbf{t}^-)$. This sign ambiguity is the
algebraic origin of the two-pose ambiguity discussed earlier.

\paragraph{Cheirality filter.}
For each candidate we transform the 3D corners and discard any pose with
$Z \leq 0$ for any corner (marker would be behind the camera --- physically
impossible).

\subsection{Block D --- Orthogonalize \texorpdfstring{$\mathbf{R}$}{R} via SVD Procrustes}

The matrix $[\mathbf{r}_1\,|\,\mathbf{r}_2\,|\,\mathbf{r}_3]$ assembled in Block~C is
\emph{approximately} a rotation matrix but not exactly: noise and numerical
errors leave the columns slightly non-orthonormal. We project to the
nearest matrix in $\SO{3}$ using the \textbf{orthogonal Procrustes}
closed-form solution.

\paragraph{Theorem (orthogonal Procrustes).}
Given any $\mathbf{M} \in \R^{3 \times 3}$, the rotation matrix
$\mathbf{R} \in \SO{3}$ closest to $\mathbf{M}$ in the Frobenius norm is
\begin{equation}
  \mathbf{M} \;=\; \mathbf{U}\,\boldsymbol{\Sigma}\,\mathbf{V}^\top
  \;\;\Longrightarrow\;\;
  \mathbf{R} \;=\; \mathbf{U}\,\mathbf{V}^\top.
\end{equation}
The singular values $\boldsymbol{\Sigma}$ are discarded; only the orthonormal
frame survives.

\paragraph{The reflection fix.}
If $\det(\mathbf{U}\mathbf{V}^\top) = -1$ the result is a reflection, not a
rotation. Flip the sign of the last column of $\mathbf{U}$ to recover
$\det = +1$ without leaving the optimal subspace:
\begin{equation}
  \text{if } \det(\mathbf{U}\mathbf{V}^\top) < 0: \;
  \mathbf{U}_{:,3} \leftarrow -\mathbf{U}_{:,3}, \;
  \mathbf{R} \leftarrow \mathbf{U}\,\mathbf{V}^\top.
\end{equation}

\subsection{Block E --- Rank candidates by reprojection error}

For each surviving candidate $(\mathbf{R}, \mathbf{t})$, compute the
\emph{reprojection error}
\begin{equation}
  E(\mathbf{R}, \mathbf{t})
  \;=\; \frac{1}{N}\sum_{i=1}^{N} \norm{\,\mathbf{p}^{(i)} - \pi(\KK(\mathbf{R}\Pmark^{(i)} + \mathbf{t}))\,}.
\end{equation}
The candidate with the smaller $E$ is reported as Solution~1 (most likely
true pose), the other as Solution~2 (the ambiguity twin). In the
noiseless case Solution~1 has $E = 0$ exactly.

\subsection{Relationship to OpenCV's implementation}

OpenCV's actual IPPE\_SQUARE follows a slightly different route:
it builds the homography using a closed-form formula specific to the
centered square (instead of generic DLT), then extracts the rotation from
the \emph{Jacobian of the homography at the marker center} via an
intermediate frame $\mathbf{R}_v$. It also applies a final iterative
refinement (Levenberg--Marquardt). Our DLT + decomposition + Procrustes
route produces \textbf{mathematically equivalent results in the noiseless
case}, and differs only by a small refinement under noise. The conceptual
flow (normalize $\to$ homography $\to$ two-fold pose decomposition $\to$
rank by reprojection error) is identical.

% ============================================================
\section{Step 7 --- Comparison with Ground Truth}
% ============================================================

Given the estimated pose $(\hat{\mathbf{R}}, \hat{\mathbf{t}})$ and the ground-truth
pose $(\Rgt, \tgt)$, we report three numbers.

\subsection{Reprojection error}

\begin{equation}
  E(\hat{\mathbf{R}}, \hat{\mathbf{t}})
  \;=\; \frac{1}{N}\sum_{i=1}^{N}
  \norm{\,\mathbf{p}^{(i)} - \pi(\KK(\hat{\mathbf{R}}\Pmark^{(i)} + \hat{\mathbf{t}}))\,}
  \quad\text{(pixels)}.
\end{equation}
Reported for each candidate solution returned by the estimator.

\subsection{Translation error}

The Euclidean distance between the estimated and true translations:
\begin{equation}
  e_t \;=\; \norm{\hat{\mathbf{t}} - \tgt} \quad\text{(meters)}.
\end{equation}

\subsection{Rotation error}

The angular distance between $\hat{\mathbf{R}}$ and $\Rgt$ is read off the
relative rotation $\mathbf{R}_{\text{err}} = \Rgt^\top\,\hat{\mathbf{R}}$. For any
rotation matrix, the trace identity gives
\begin{equation}
  \mathrm{trace}(\mathbf{R}) \;=\; 1 + 2\cos\theta,
\end{equation}
so
\begin{equation}
  \boxed{\quad
    e_R \;=\; \arccos\!\left(\frac{\mathrm{trace}(\Rgt^\top \hat{\mathbf{R}}) - 1}{2}\right)
  \quad\text{(radians; convert to degrees as needed)}.
  \quad}
  \label{eq:rot-error}
\end{equation}
If $\hat{\mathbf{R}} = \Rgt$, then $\mathbf{R}_{\text{err}} = \mathbf{I}$, trace is $3$, and
$e_R = 0$. The output lives in $[0, \pi]$ regardless of the axis of
disagreement, giving a single intuitive ``how off is the rotation'' scalar.

\paragraph{Numerical guard.}
Floating-point round-off can push the argument of $\arccos$ slightly
outside $[-1, 1]$, producing \verb|NaN|. The code clamps the argument
to $[-1, 1]$ before the $\arccos$ call.

% ============================================================
\section*{Summary Table}
% ============================================================

\begin{center}
\begin{tabular}{|c|l|l|}
\hline
\textbf{Step} & \textbf{What it produces} & \textbf{Code location} \\
\hline
1   & Marker plane ($Z=0$)                       & implicit in \texttt{marker.py} \\
2   & Square corners $\Pmark^{(i)}$              & \texttt{square\_object\_points(L)} \\
3   & Camera intrinsics $\KK$                    & \texttt{make\_default\_camera()} \\
4   & Pose $(\Rgt, \tgt, \wgt)$                  & \texttt{make\_ground\_truth\_pose()} \\
5a  & Camera-frame points $\Pcam^{(i)}$          & \texttt{transform\_points(...)} \\
5b  & Image points $\mathbf{p}^{(i)}$            & \texttt{camera.project(...)} \\
6   & Estimated pose $(\hat{\mathbf{R}}, \hat{\mathbf{t}})$ & \texttt{ippe\_square(...)} \\
7   & Errors $E, e_t, e_R$                       & \texttt{pipeline.py} (reporting) \\
\hline
\end{tabular}
\end{center}

\end{document}
